\documentclass[../main.tex]{subfiles}
\ifSubfilesClassLoaded{\setcounter{chapter}{16}}{}
\begin{document}

\chapter{Presentasi Hasil Proyek dan Penyerahan Akhir}

\begin{subcpmk}
  \item \textbf{Sub-CPMK-P10:} Memresentasikan masalah, metode, hasil, dan rekomendasi dalam waktu 5--10 menit; menyerahkan notebook final (RPS Mg.\ 16).
\end{subcpmk}

\SesiRPS{16}{P10}{$\sim$170 menit}{CPMK-P4}

\noindent Presentasi menutup siklus OBE: bukti angka diubah menjadi narasi keputusan. Penilaian mengutamakan kejelasan argumen dan pengakuan keterbatasan, bukan hanya metrik tertinggi.

\section{Struktur presentasi}

Alokasi waktu di bawah dapat disesuaikan dengan rubrik dosen; gunakan timer saat latihan agar bagian rekomendasi tidak terpotong.

\begin{enumerate}[leftmargin=*]
  \item Masalah dan motivasi (30--60 dtk).
  \item Data dan prapemrosesan ringkas (60 dtk).
  \item Model dan evaluasi (2--3 mnt).
  \item Temuan utama dan rekomendasi (1--2 mnt).
  \item Keterbatasan dan etika data (30 dtk).
\end{enumerate}

\begin{table}[htbp]
\centering
\caption{Alokasi waktu presentasi (panduan 8--10 menit).}
\label{tab:waktu-presentasi}
\begin{tabular}{@{}p{4.5cm}cc@{}}
\toprule
\textbf{Bagian} & \textbf{Menit min} & \textbf{Menit maks} \\
\midrule
Pembukaan & 0.5 & 1 \\
Data \& persiapan & 1 & 1.5 \\
Model \& metrik & 2 & 3 \\
Temuan \& rekomendasi & 1.5 & 2.5 \\
Penutup (risiko/etika) & 0.5 & 1 \\
\bottomrule
\end{tabular}
\end{table}

\section{Checklist penyerahan}

\begin{itemize}[leftmargin=*]
  \item Notebook final: \textit{Restart \& Run All} lancar.
  \item Berkas dinamai sesuai konvensi Lampiran.
  \item Slide PDF/PNG (jika diminta) + laporan ringkas.
\end{itemize}

\begin{figure}[htbp]
\centering
\begin{tikzpicture}[node distance=6mm]
  \node[dsboxwide, fill=purple!8] (m) {Bab PDF\\modul};
  \node[dsboxwide, fill=purple!8, right=of m] (n) {Notebook\\\texttt{.ipynb}};
  \node[dsboxwide, fill=purple!8, right=of n] (r) {Run All\\+ narasi};
  \draw[dsarrow] (m) -- (n);
  \draw[dsarrow] (n) -- (r);
\end{tikzpicture}
\caption{Alur belajar disarankan: baca tujuan sesi di PDF, lalu implementasikan dan dokumentasikan di notebook hingga seluruh sel dapat dijalankan ulang dari awal.}
\end{figure}


\section{\texorpdfstring{Contoh slide (Markdown $\rightarrow$ export)}{Contoh slide Markdown ke export}}

Mahasiswa dapat memakai Jupyter + Rise, Marp, atau PowerPoint; contoh poin di Markdown:

\MulaiPenjelasanKode Struktur heading Markdown memetakan satu ide per slide: konteks data, bukti model, aksi bisnis---memudahkan audiens non-teknis mengikuti alur. \SelesaiPenjelasanKode

\begin{verbatim}
# Slide judul
## Data: sumber, N baris, target
## Model terpilih + metrik utama
## Rekomendasi bisnis (3 bullet)
\end{verbatim}

\section{Contoh kode (demo cepat di kelas)}

\MulaiPenjelasanKode Menunjukkan pola \texttt{joblib.load} untuk membuka artefak model yang disimpan minggu lalu; baris aktif di bawah hanya placeholder agar sel tetap jalan tanpa berkas \texttt{.joblib}. \SelesaiPenjelasanKode

\begin{lstlisting}[caption=Muat ulang model/figur final]
# from joblib import dump, load
# load("model_final.joblib")  # jika disimpan
print("Notebook final siap ditunjukkan")  # ganti dengan load/prediksi nyata bila ada artefak
\end{lstlisting}

\MulaiPenjelasanKode \texttt{os.path.getsize} memberi ukuran byte; dibagi 1024 untuk KB agar Anda memastikan unggahan di bawah batas portal pengumpulan. \SelesaiPenjelasanKode

\begin{lstlisting}[caption=Cek ukuran berkas sebelum unggah]
import os
for f in ["proyek_final.ipynb", "laporan.pdf"]:  # sesuaikan daftar nama berkas penyerahan Anda
    if os.path.exists(f):  # cegah error jika berkas belum dibuat
        print(f, os.path.getsize(f) / 1024, "KB")  # ukuran dalam kilobyte (perkiraan)
\end{lstlisting}

\section{Referensi}

\begin{itemize}[leftmargin=*]
  \item INRIA sklearn MOOC. \url{https://inria.github.io/scikit-learn-mooc/}
\end{itemize}

\begin{asesmen}
\textbf{Presentasi \& proyek akhir} sesuai rubrik RPS Bagian 8.
\end{asesmen}

\begin{rangkuman}
Minggu 16 menutup siklus komunikasi hasil dan bukti capaian integratif.

\textbf{Kata kunci:} presentasi, penyerahan, refleksi
\end{rangkuman}

\end{document}
